\chapter{Linux systems and their resilience needs}
The family of linux operating systems, based on the Linux kernel, has proven to be very versatile and performant in the most diverse fields of computation. Embedded Systems, cloud systems, desktops or real-time systems are often operated by a Linux Kernel.\\
Supporting these diverse systems and their complex requirements is no easy feat, but they have one common demand: Resilience

Information systems often serve critical purposes, and must be resilient towards faults and provide high uptimes. To achieve that, literature provides many resilience patterns, which help to mitigate and react to faults or failures. 

The Linux Kernel is certainly a trusted piece of software, and studying its implemented resilience patterns is the topic of this work. Understanding the decisions of the Kernel maintainers is valuable for developers who seek to improve their own systems. Especially the open source code and public discussions allow for a deep insight, as opposed to similar systems of Apple or Microsoft.

The implemented resilience patterns and those that could have been used, but did not provide enough benefit, can help operating system developers help improve their own systems.
